\section{Running the project}
\label{app:running}

The code, the maps, the design log and this report are in one repository.
It needs Webots R2025a and \texttt{uv}, which installs Python and every
package at the versions in the lock file. On Windows, start Webots and the
tools from PowerShell.

{\small
\begin{verbatim}
uv sync                                  # the environment
uv run make_world.py dead_end -n 20 -j 2 -s 4
                     # a world (map, drones, leaders, seed): open
                     # worlds/dead_end.wbt in Webots
uv run tools/sweep.py --map dead_end -n 10 20 30 -j 0 1 2 4 \
    --seeds 0 1 2 3 4 5 6 7 8 9 --duration 83 --model hybrid --plan 1
uv run tools/analyse.py all              # every number and figure
uv run pytest                            # the tests
cd report; tectonic main.tex             # this report
\end{verbatim}}

Any setting of any layer can be changed for one run by its name, for example
\texttt{-{}-model couzin -{}-velocity 0.5 -{}-plan 1}. The design log,
\texttt{docs/design-log.md}, records every decision with its reason and the
alternatives turned down.

\section{Videos}
\label{app:videos}

The videos are in one shared folder,
\begin{center}
  \small\expandafter\url\expandafter{\videos}
\end{center}
each named after the figure it belongs to. They were filmed with
\texttt{tools/record.py}, which re-runs a sweep's world from its seed; the
counter in the corner is the playback speed. All have 20 drones and the
leaders planning.
\begin{description}
  \item[\Cref{fig:video-forest}:] the hybrid on \texttt{l\_hall}, 2 leaders,
    seed 0; the camera follows the flock, tilted.
  \item[\Cref{fig:video-reynolds}:] Reynolds on \texttt{dead\_end}, 2
    leaders, seed 2; the whole map from above.
  \item[\Cref{fig:video-couzin}:] Couzin on \texttt{dead\_end}, seed 2, one
    video with 1 leader and one with 4; the whole map from above.
\end{description}
